\section{Group-quotient localizers}

Write $\mathcal C_{\mathfrak a}$ for the coefficient groups in LM8,
B.6: respectively all abelian groups, all metabelian groups, groups
nilpotent of class at most $c$, finite $\ell$-groups, the groups
$U(\mathbb Q)$ for finite-dimensional unipotent algebraic groups
$U/\mathbb Q$, and all virtually abelian discrete groups. Throughout
this section an $\mathfrak a$-equivalence of groups means a homomorphism
$v:G'\to G$ such that
\[
 \Hom(G,H)\longrightarrow\Hom(G',H)
       \quad\hbox{is bijective for every }H\in\mathcal C_{\mathfrak a}.
                                                        \tag{H5.0.1}
\]
A map of spaces is in $W_{\mathfrak a}$ precisely when it is bijective
on components and its component group homomorphisms satisfy
(H5.0.1). This formulation includes the continuous or algebraic
universal properties of the two completions; it never replaces a
completion by its underlying abstract group.

\begin{claim}{H5.1}{Proved}{The groupoid orthogonality used here}
For a homomorphism $v:G'\to G$ and a group $H$, condition
$\Hom(G,H)\cong\Hom(G',H)$ is equivalent to
$\Fun(\mathbf BG,\mathbf BH)\to\Fun(\mathbf BG',\mathbf BH)$
being an equivalence. Consequently (H5.0.1), with the separate
component condition, is exactly the definition of LM8, B.6.
\end{claim}
\begin{proof}
The functor groupoids have homomorphisms as objects and conjugating
elements of $H$ as arrows. If restriction is bijective on objects,
an element conjugating two restrictions conjugates their unique
extensions, proving full faithfulness. Conversely, essential
surjectivity lets one conjugate an extension so that its restriction
equals a prescribed homomorphism exactly. If two exact extensions
have equal restrictions, fullness lifts the identity conjugating
element; that same element is the identity in $H$, so the extensions
are equal. This is also \citet[Theorem 1.3]{EH5CGT2006}.

The checked calculation in INDLOCAL, B.6 is the unpointed version:
its product of sets $\Hom(-,H)/H$ is precisely the set of
components of this functor groupoid. Its B.3 follows from the
canonical first truncation and, for $W_0$, the component reflection.
For B.4 the relevant classical class consists of group maps that
are isomorphisms on $H_1(-;\mathbb Z)$ and epimorphisms on
$H_2(-;\mathbb Z)$. Bousfield's group-localization theorem and
realization lemma are \citet[Theorems 5.2, 5.5 and Lemma 6.1]{EH5Bousfield1975}.
The realization lemma attaches $1$- and $2$-cells to realize the
group map, then $3$-cells to kill the free, spherical relative
$H_2$. Thus a homology-invariant groupoid representable is
orthogonal to the group-localization units; conversely the Hopf
exact sequence puts the fundamental-group map of every homology
equivalence in the indicated group class. This proves the asserted
description of the local groups. It does not identify an arbitrary
localizer with the maps detected by its small representables.
\end{proof}

\begin{claim}{H5.2}{Refuted}{Reverse cofibration and homology containment are different}
None of the six localizers $W_{\mathfrak a}$ satisfies Loc~(6 b).
Nevertheless
\[
 W_{\mathbb Z}\subset W_{\mathfrak a}
 \quad\Longleftrightarrow\quad
 \mathfrak a\in
 \{\mathrm{ab},\mathrm{nil}_c,\mathrm{pro}\mbox{-}\ell,
                         \mathrm{unip}_{\mathbb Q}\}.
                                                        \tag{H5.2.1}
\]
Thus the asserted equivalence in Part H.5(ii) is false.
\end{claim}
\begin{proof}
First consider the four-point poset $P$ with two minimal elements,
two maximal elements, and all four possible strict comparisons
between these levels. Its nerve is a circle. Add two incomparable
maximal cone vertices $u,v$, both above all of $P$, and call the
result $X$. Put $U=P\cup\{u\}$ and $V=P\cup\{v\}$. These are
sieves in $X$ with $X=U\cup V$ and $U\cap V=P$.
The nerves of $U,V$ are contractible and $BX\simeq S^2$.
They are therefore $W_{\mathfrak a}$-aspherical. The circle $BP$
is not: every coefficient class contains a nontrivial group and
$\Hom(\mathbb Z,H)$ is then nontrivial. This is the reverse
cofibration counterexample. The complete poset tables and integral
Smith forms occur in
\texttt{code/outputs/EH5\_group\_witnesses.json}; the Betti numbers
are $(1,1)$ for $P$, $(1)$ for each cone, and $(1,0,1)$ for $X$.

Every nilpotent group $H$ has $BH$ integral-homology local.
Here is a direct verification of the fact needed in this argument.
All $K(A,n)$, $n\geq1$, are integral-homology local by the universal
coefficient theorem and the formula for their mapping-space
homotopy groups. A central extension
$1\to A\to H\to Q\to1$ expresses $BH$ as the homotopy fiber
of its extension-class map $BQ\to K(A,2)$. Induction on a finite
central series proves the assertion, because local spaces are
closed under homotopy limits. All coefficient groups in the four
positive cases are nilpotent: finite $\ell$-groups have a finite
central series, and a finite-dimensional unipotent group in
characteristic zero is nilpotent. Hence every integral homology
equivalence is inverted by each coefficient target, giving the
four inclusions.

For both negative cases use $F=F(x,y)$ and the endomorphism
\[
 \alpha(x)=x,\qquad
 \alpha(y)=yx^{-1}y^{-1}xy.                              \tag{H5.2.2}
\]
Its exponent-sum matrix is the identity. Since $BF$ is a graph,
$B\alpha$ is an integral homology equivalence.
Over $\Lambda=\mathbb Z[X^{\pm1},Y^{\pm1}]$ its abelianized Fox
Jacobian, with columns corresponding to the two images, is
\[
 J(\alpha)=
 \begin{pmatrix}
 1&(1-Y)/X\\
 0&2-X^{-1}
 \end{pmatrix}.                                         \tag{H5.2.3}
\]
It cannot induce an automorphism of $F/F''$. Indeed Fox
derivatives followed by abelianization vanish on $F''$.
An inverse on $F/F''$ could be lifted on the two free generators;
the Fox chain rule, and the identity on abelianization, would then
make (H5.2.3) invertible over $\Lambda$. Its determinant
$2-X^{-1}$ is not a Laurent unit. Thus $B\alpha\notin W_{\mathrm{meta}}$.

For the virtually abelian case take the allowed finite target
$S_3$ and the homomorphism $\rho(x)=(12)$, $\rho(y)=(13)$.
Then $\rho(F)=S_3$, whereas
$(\rho\alpha)(x)=(\rho\alpha)(y)=(12)$.
The automorphism groups of these objects in the two functor
groupoids are respectively $C_{S_3}(S_3)=1$ and
$C_{S_3}(\langle(12)\rangle)\cong C_2$. Restriction is not fully
faithful, so $B\alpha\notin W_{\mathrm V}$.

The program \texttt{code/EH5\_group\_witnesses.py} checks the Fox
matrix and both centralizers. It also realizes (H5.2.2) by a
functor between finite incidence categories: subdivide the
$y$-circle into the five oriented edges in the displayed word,
and leave the $x$-circle unsubdivided. The target incidence
category has $3$ objects and $7$ arrows; the source has $11$
objects and $23$ arrows. All objects, face arrows, compositions,
and the functor table are in its output. Thus the negative
containment witnesses can be taken in finite $\Cat$.
\end{proof}

\begin{claim}{H5.3}{Proved}{The three algebraic quotient cases are accessible}
For abelianization, metabelianization, and the quotient of
nilpotence class $c$, $W_{\mathfrak a}$ is accessible.
A generating set consists of $S^2\to *$ and the map
$R_r\to K_w$ that attaches a $2$-cell to the rose on $r$
circles along the word
\[
 \begin{array}{c|c|c}
 \mathfrak a&r&w\\ \hline
 \mathrm{ab}&2&[x_1,x_2]\\
 \mathrm{meta}&4&[[x_1,x_2],[x_3,x_4]]\\
 \mathrm{nil}_c&c+1&[x_1,[x_2,\ldots,[x_c,x_{c+1}]\ldots]] .
 \end{array}                                            \tag{H5.3.1}
\]
Here and below a space map in a generating set means its
category-of-simplices representative.
\end{claim}
\begin{proof}
Locality for $S^2\to *$ says exactly that every component is a
$1$-type. For a $1$-type $BH$, locality for $R_r\to K_w$ says
that every $r$-tuple in $H$ satisfies $w=1$. There is no extra
condition on arrows in the mapping groupoids: the group map
$F_r\to F_r/\langle\!\langle w\rangle\!\rangle$ is surjective,
so centralizers of an extension and its restriction are equal.
The three laws define exactly the three indicated varieties.
The resulting local spaces are arbitrary disjoint unions of
their classifying spaces, and their detected equivalences are
(H5.0.1). The group reflectors are respectively $G/G'$,
$G/G''$, and $G/\Gamma_{c+1}G$.
\end{proof}

\begin{claim}{H5.4}{Proved}{Accessibility for finite and virtually abelian tests}
The pro-$\ell$ and virtually abelian localizers are accessible.
In each case a generating set consists of $S^2\to *$
and all its group equivalences $BG'\to BG$ with
$|G'|,|G|\leq\gamma$, for a single
$\mathcal U$-cardinal $\gamma$ constructed below.
These maps are then represented by categories of simplices.
\end{claim}
\begin{proof}
We first prove a bounded-generation lemma for a compact
Hausdorff group $C$. Let $v:G'\to G$ induce a bijection
$\Hom(G,C)\to\Hom(G',C)$, and start with subgroups
$A\subseteq G'$, $B\subseteq G$, $v(A)\subseteq B$.
There is a countable enlargement procedure with the following
three requirements.

First, for each homomorphism $A\to C$ not extending to $G'$,
adjoin finitely many elements of $G'$ on whose generated
extension it already fails to extend. Such a finite
obstruction exists because homomorphisms into $C$ are the
closed solution sets of the multiplication equations in
the compact product $C^{G'}$. Do the same for $B\subseteq G$.
Second, enlarge $B$ to contain the image of the new $A$.
Third, for each $b\in B$ and distinct $c,d\in C$, consider
the compact space of pairs $\phi,\psi:G\to C$ with
$\phi(b)=c$, $\psi(b)=d$. No such pair agrees on all
of $v(G')$, by injectivity of restriction. The closed
equalities $\phi(v(a))=\psi(v(a))$, $a\in G'$, therefore
have empty intersection in that compact space. Finitely
many of them already have empty intersection. Adjoin
these finitely many $a$ to $A$.

Repeat these operations countably, and write
$A_\infty,B_\infty$ for the unions. Every map
$A_\infty\to C$ extends to $G'$: each restriction to
a stage extends, and the compact spaces of global
extensions form a nested family with nonempty
intersection. Similarly every map $B_\infty\to C$
extends to $G$. Surjectivity of
$\Hom(B_\infty,C)\to\Hom(A_\infty,C)$ now follows
from the global equivalence $v$. For injectivity,
extend two maps on $B_\infty$ globally. If they
agree on $A_\infty$ but differ at some $b$ with
values $c,d$, the third enlargement at a stage
containing $b$ gives a contradiction. Thus the
restricted map is a $C$-equivalence.

This enlargement has a bound depending only on
$|A|+|B|+|C|+\aleph_0$. At a stage of size at most
$\lambda\geq |C|+\aleph_0$, there are at most
$2^\lambda$ maps and pairs to consider; adjoining
the finite witnesses and generating subgroups
has that same bound. Iterating the power-set bound
countably gives a cardinal $b(\lambda)\in\mathcal U$.
To obtain a directed presentation with one uniform
bound, perform the construction coherently for
every finite subset $S$ of $G'\amalg G$. Induct on
$|S|$, requiring its chosen enlargement to contain
the finitely many enlargements for proper subsets
of $S$. If $\lambda_0=|C|+\aleph_0$ and
$\lambda_{r+1}=b(\lambda_r)$, then all chosen
submaps have size at most
$\gamma=\sup_{r<\omega}\lambda_r\in\mathcal U$.
They form a filtered diagram whose colimit is $v$.
Nerves commute with filtered colimits of groups,
and filtered colimits of simplicial sets compute
their homotopy colimits. Hence every $C$-equivalence
is generated by the indicated bounded maps.
Together with $S^2\to *$ this proves the lemma
for the associated fundamental localizer.

For pro-$\ell$ use the compact group which is
the product of representatives of all finite
$\ell$-groups. Testing its Hom functor is
equivalent to testing all its factors, so the
lemma applies.

For the virtually abelian case we verify that compact targets
of bounded size suffice. Let $\mathcal K$ be the set of all
compact Lie groups that are finite extensions of tori, viewed
as abstract groups. Every virtually abelian group belongs to
the closure of $\mathcal K$ under small limits. To see this,
write
$1\to A\to H\to F\to1$ with $A$ abelian and $F$ finite.
The normal abelian subgroup can be obtained by intersecting
the finitely many conjugates of an abelian subgroup of finite
index. Embed the $F$-module $A$ equivariantly into
$D=(\mathbb R/\mathbb Z)^I$ with $F$ permuting coordinates:
use all characters of $A$ into the injective cogenerator
$\mathbb R/\mathbb Z$, together with their $F$-translates.
Push the extension out along $A\hookrightarrow D$ to obtain
$H_D$. Its underlying set is $D\times F$, with multiplication
given by the original action and a $2$-cocycle. With the
product topology it is a compact group.

The group $H_D$ is the inverse limit of its quotients formed
by projecting $D$ onto finite $F$-stable coordinate sets.
Each quotient is a finite extension of a finite-dimensional
torus, and is therefore in $\mathcal K$. The quotient
$H_D/A$ splits as $(D/A)\rtimes F$, since the defining
cocycle has values in $A$. Embed the $F$-module $D/A$ into a
second permutation product of circles $D'$. Then $H$ is the
equalizer of the two homomorphisms
\[
 H_D\rightrightarrows D'\rtimes F:
 \qquad
 (d,f)\longmapsto(\overline d,f),\quad
 (d,f)\longmapsto(0,f).
                                                        \tag{H5.4.1}
\]
The target $D'\rtimes F$ is again an inverse limit of members
of $\mathcal K$. This proves the limit assertion. Since
orthogonality of group homomorphisms is preserved by limits
of target groups, testing $\mathcal K$ is equivalent to
testing all virtually abelian groups. Conversely every
member of $\mathcal K$ is virtually abelian as an abstract
group, so no additional test has been introduced.
Now take the compact product of representatives
of $\mathcal K$ and apply the proved lemma.
All its factors have cardinality at most the
continuum, so the product and the resulting
bound remain in $\mathcal U$.

This direct group argument avoids a delicate
point in the opened proof of
\citet[Theorem 19]{EH5Przezdziecki2002}.
The asserted equivalence between the homotopy
category of arrows and arrows in the homotopy
category, following its Definition 10, is
false in general: it forgets the homotopy
specified in a homotopy-commutative square.
Our argument uses only closed sets of actual
group homomorphisms, and does not require
that equivalence or the affected arrow
category reduction.
\end{proof}

\begin{claim}{H5.5}{Proved}{None of the six cases is a stable inverse image}
No $W_{\mathfrak a}$ in the list is of the form $R(\mathcal L)$.
In fact $T(W_{\mathfrak a})=\Ho(\mathrm{Sp})$, so its stable
closure $R(T(W_{\mathfrak a}))$ is all maps.
\end{claim}
\begin{proof}
Every $W_{\mathfrak a}$ contains $S^2\to *$, whose unreduced
homotopy cofiber is $S^3$. A stable localizing subcategory
containing its suspension spectrum contains the sphere and
hence every spectrum. But $S^1\to *$ is not in any of the
six classes, as checked in H5.2. Thus none is the all-map
class, proving both assertions.
\end{proof}

\begin{claim}{H5.6}{Proved}{All intrinsic fibrations: an exact group-theoretic test}
Let $\mathcal C=\mathcal C_{\mathfrak a}$ be any of the six
coefficient classes. For a functor $p:E\to B$, put
$P=\Pi_1E$, $Q=\Pi_1B$, and let $F_b$ denote the homotopy
fiber of $BE\to BB$. Then $p\in\Fib_{W_{\mathfrak a}}$
if and only if all three conditions below hold.
\begin{enumerate}
\item For every $\sigma:[n]\to B$ and vertex $i$,
$E_{\sigma(i)}\to E\times_B[n]$ is in $W_{\mathfrak a}$.
\item For every $b$, the natural map
\[
 F_b\longrightarrow
       \fib_b(BP\longrightarrow BQ)                     \tag{H5.6.1}
\]
is in $W_{\mathfrak a}$. Equivalently, every homomorphism
$\pi_1(F_b,x)\to H$, $H\in\mathcal C$, kills the image of
the boundary $\pi_2(BB,b)\to\pi_1(F_b,x)$.
\item The following explicit group test holds on each
component of $Q$. Write that component as $\mathbf BH$
and the components of $P$ above it as
$\mathbf BG_j\xrightarrow{\alpha_j}\mathbf BH$.
For every $\mathfrak a$-equivalence $\beta:K'\to K$
and every $\theta:K\to H$, the map
\[
 \coprod_j\alpha_j(G_j)\backslash H/\theta\beta(K')
 \longrightarrow
 \coprod_j\alpha_j(G_j)\backslash H/\theta(K)              \tag{H5.6.2}
\]
is bijective, and for each representative $h\in H$ the map
\[
 \{(g,k'):\alpha_j(g)=h\theta\beta(k')h^{-1}\}
 \longrightarrow
 \{(g,k):\alpha_j(g)=h\theta(k)h^{-1}\}                    \tag{H5.6.3}
\]
is an $\mathfrak a$-equivalence of groups.
\end{enumerate}
The phrase $\mathfrak a$-equivalence in (3) means the explicit
Hom test (H5.0.1). Thus this is an equality describing
$\Fib_{W_{\mathfrak a}}$ on all functors, with no restriction
to coverings or to finite categories.
\end{claim}
\begin{proof}
We first record why condition (1) handles strictness.
For a simplex $\sigma:\Delta[n]\to K$ over $NB$, set
$P_\sigma=NE\times_{NB}\Delta[n]$. The augmentation
$\hocolim_{\Delta/K}P_\sigma\to NE\times_{NB}K$ is a weak
equivalence, by the simplex resolution used in E.3.
For each target $BH$, (1) says that the contravariant
diagram $\Map(P_\sigma,BH)$ has equivalences as transition
maps. Its homotopy limit is the space of sections of a
locally constant diagram over $|K|$. A weak equivalence
of bases induces an equivalence of these section spaces.
Thus strict pullback preserves weak equivalences in
$W_{\mathfrak a}$. Pulling back a path fibration from a
Kan replacement of $NB$ identifies strict and homotopy
fibers in $W_{\mathfrak a}$. A comparison with an ordinary
fibration replacement, simplex by simplex and then by
the same homotopy-limit formula, identifies every strict
pullback with its derived pullback in $W_{\mathfrak a}$.
The last-vertex comparison for categories of simplices
therefore transfers the intrinsic property between
categorical bases and arbitrary spaces over $BB$.

For a space map $e:X\to Y$, first truncation gives a
comparison
\[
 \Pi_1(X\times_Y^h Z)\longrightarrow
   \Pi_1X\times_{\Pi_1Y}^h\Pi_1Z.                        \tag{H5.6.4}
\]
It is bijective on components and surjective on each
isotropy group. Its kernel is the image of the appropriate
boundary from $\pi_2Y$; this follows either from the
homotopy exact sequence, or from the path model of both
pullbacks. If (2) holds, every representation of the
isotropy of the left side in an allowed $H$ kills this
kernel: restrict it first to the corresponding fiber
$\pi_1F_b$. Thus (H5.6.4) is a $W_{\mathfrak a}$-equivalence
for every $Z\to Y$.

The homotopy pullback of
$\coprod_j\mathbf BG_j\to\mathbf BH$ along
$\mathbf BK\to\mathbf BH$ has exactly the component set
in the target of (H5.6.2), and its isotropy at the
representative $h$ is the target group in (H5.6.3).
Consequently (3) says precisely that these groupoid
pullbacks preserve all $\mathfrak a$-equivalences.
Componentwise decomposition covers disconnected
groupoids as well. Conditions (1)--(3), (H5.6.4), and
two out of three now prove sufficiency.

For necessity, an intrinsic $p$ satisfies (1) by pulling
back $[0]\to[n]$. Hence it has the derived preservation
property just proved. Pull it back over the universal
cover of a component of $BB$. The inclusion of its
basepoint is a $W_1$-equivalence, so its pullback gives
(2), by the homotopy exact sequence.
Finally any groupoid map $R\to Q$ is realized over $BB$
by $BB\times_{BQ}^h BR\to BB$; its first truncation is
$R$. A $W_{\mathfrak a}$-equivalence $R'\to R$ gives a
$W_{\mathfrak a}$-equivalence of these base spaces.
Their total pullbacks are $BE\times_{BQ}^h BR'$ and
$BE\times_{BQ}^h BR$, whose fundamental groupoids are
the corresponding groupoid pullbacks, because $BQ$
is a $1$-type. Intrinsic preservation yields exactly
(H5.6.2)--(H5.6.3).
\end{proof}

\begin{claim}{H5.7}{Refuted}{Local groups alone do not classify all larger localizers}
The assignment
$W\mapsto\{G:BG\text{ is }W\text{-local}\}$ is not injective,
even on accessible fundamental localizers containing $W_1$.
Three different such localizers have exactly the trivial
group as a local group.
\end{claim}
\begin{proof}
The first is $W_0$, the maps bijective on components,
generated by $S^1\to *$. Its local spaces are discrete
sets. The second is $W_{-1}$, the maps preserving
emptiness versus nonemptiness, generated by $S^0\to *$.
Indeed locality means that the diagonal
$Z\to Z\times Z$ is an equivalence; its solutions are
the empty space and the contractible spaces. Maps
detected by them are exactly those preserving
emptiness. The third is the all-map localizer,
generated by $\varnothing\to *$, with contractible
local spaces. In each case $BG$ is local exactly when
$G=1$, whereas
\[
 W_1\subset W_0\subsetneq W_{-1}\subsetneq\mathrm{Fl}(\Cat).
                                                        \tag{H5.7.1}
\]
For completeness $W_{-1}$ satisfies relative Theorem A
directly. In a triangle $A\to B\to C$, suppose every
$A/c\to B/c$ preserves emptiness. If $B$ is nonempty,
choose $b\in B$ and let $c$ be its image; then
$B/c$ is nonempty, hence so are $A/c$ and $A$.
If $B$ is empty then so is $A$. Weak saturation and
the terminal-object axiom are immediate. This checks
the convention of fundamental localizer in this paper.
\end{proof}

\begin{claim}{H5.8}{Proved}{The accessible component-preserving correspondence}
On the interval $W_1\subseteq W\subseteq W_0$, accessible
fundamental localizers correspond exactly to small
orthogonality classes of groups. Equivalently they
correspond to full reflective, accessibly embedded
subcategories of groups. A localizer associated to
$\mathcal D$ has local spaces the arbitrary disjoint
unions of $BG$ with $G\in\mathcal D$, and its maps are
component bijections inducing isomorphisms after the
group reflector.
The local groups are closed under all small limits
and retracts, and under $\lambda$-filtered colimits
for some regular $\lambda$. They need not be closed
under extensions.
\end{claim}
\begin{proof}
Localizing at $W_1$ first restricts to groupoids.
The component condition ensures that every discrete
set is local and that localization preserves
components. Hence localization is groupwise, as
also proved in \citet[Theorems 2.6--2.7 and 3.3]{EH5CGT2006}.
A small set of groupoid generators may be replaced
by the set of its component group homomorphisms.
H5.1 then identifies its local isotropy groups with
their ordinary orthogonality class.
Conversely a small set $\{G_i\to H_i\}$ of group
maps, together with $S^2\to *$, generates the stated
fundamental localizer. The local group reflector
exists by the ordinary small-orthogonality
construction in the locally presentable category
of groups; \citet[Theorem 3.1]{EH5CGT2006} identifies
its inverted maps with the mapping-groupoid tests.

Limits and retracts preserve orthogonality.
Choose a regular $\lambda$ for which all generator
groups are $\lambda$-presentable; their Hom functors
commute with $\lambda$-filtered colimits, proving
the asserted closure. Conversely a reflective
accessibly embedded subcategory is a small
orthogonality class: choose a sufficiently large
regular presentability rank and take the reflector
units on its set of presentable groups. Expressing
an arbitrary group as the corresponding filtered
colimit proves orthogonality to every unit.
Extension closure fails already for the abelian
local groups: the integral Heisenberg group is a
nonabelian central extension of $\mathbb Z^2$
by $\mathbb Z$.

In this correspondence the first three cases
occupy their indicated varieties. In the remaining
cases the exact class is
$(\,{}^\perp\mathcal C_{\mathfrak a}\,)^\perp$,
the orthogonality closure of the stipulated
coefficient groups. It is essential to retain this
closure rather than identify it with underlying
pro-completions. H5.7 explains the additional
component information required outside the
displayed interval.
\end{proof}

\begin{claim}{H5.9}{Proved}{The rational unipotent case depends on set theory}
For the rational unipotent completion exactly as defined in LM8 B.6,
with tests $U(\mathbb Q)$ for finite-dimensional unipotent
$\mathbb Q$-groups, the following statements hold.
If the fixed universe $\mathcal U$ contains no measurable cardinal,
then $W_{\mathrm{unip}}$ is not accessible. If Vop\v enka's principle
holds for $\mathcal U$, then $W_{\mathrm{unip}}$ is accessible.
In the latter case a generating set consists of $S^2\to *$ and
all homomorphisms between groups of cardinality less than some
regular $\kappa\in\mathcal U$ which induce bijections on
$\Hom(-,U(\mathbb Q))$ for every finite-dimensional unipotent
$\mathbb Q$-group $U$, realized by classifying-space maps and
categories of simplices.
Thus an unconditional affirmative answer to H.5(i) for this case
cannot be deduced in ZFC with the stated universes; an unconditional
negative answer is incompatible with the indicated Vop\v enka
hypothesis. The resulting independence assertion is relative to
the consistency of the corresponding universe and large-cardinal
hypotheses.
\end{claim}
\begin{proof}
We give the obstruction algebra and the smallness argument in full.
For a commutative unital $\mathbb Q$-algebra $R$, let
\[
 \mathcal H(R)=R^3,\qquad
 (a,b,c)(a',b',c')=(a+a',b+b',c+c'+ab')
                                                               \tag{H5.9.1}
\]
be its Heisenberg group. Its center and commutator subgroup are
both $C(R)=\{(0,0,c):c\in R\}$, since
$[(a,0,0),(0,b,0)]=(0,0,ab)$ and every $c$ is $c\cdot1$.
It is uniquely divisible and nilpotent of class two.
In particular $\mathcal H(\mathbb Q)$ is an allowed
finite-dimensional unipotent test group.

We first establish an elementary detection lemma. If $R$ has
no nonzero finite-dimensional unital $\mathbb Q$-algebra quotient,
then every homomorphism
\[
             \rho:\mathcal H(R)\longrightarrow U(\mathbb Q)
                                                               \tag{H5.9.2}
\]
into an allowed target kills $C(R)$. Indeed use exponential and
logarithm, which are finite polynomial expressions for unipotent
matrices in characteristic zero; their inverse and multiplication
properties are \citet[Propositions 14.32 and 14.36]{EH5Milne2022}.
The Zariski closure of the image still has class at most two,
since the corresponding iterated-commutator identity is a polynomial
identity and holds on its dense subgroup. A group homomorphism respects
unique roots. On a class-two image, addition in its Lie algebra
is recovered by multiplying group elements and correcting by the
half-commutator. Hence the three maps
\[
 A(r)=\log\rho(r,0,0),\quad B(s)=\log\rho(0,s,0),\quad
 C(t)=\log\rho(0,0,t)
\]
are $\mathbb Q$-linear and satisfy
$[A(r),B(s)]=C(rs)$.
If $C$ is nonzero, choose a linear functional $\ell$ on the
finite-dimensional Lie algebra of $U$ such that
$f=\ell C:R\to\mathbb Q$ is nonzero. The bilinear form
$(r,s)\mapsto f(rs)$ has finite rank, since it factors through
$A:R\to\operatorname{Lie}U$ in its first variable. Therefore
\[
              I_f=\{r\in R:f(rs)=0\text{ for all }s\in R\}
                                                               \tag{H5.9.3}
\]
has finite codimension. It is an ideal: if $r\in I_f$ and
$a\in R$, then $f((ar)s)=f(r(as))=0$. It is proper because
$f\ne0$, so $1\notin I_f$. The quotient $R/I_f$ contradicts
the hypothesis. This proves the detection lemma, for targets
of every nilpotence class.

Let $\kappa$ be any infinite regular cardinal in $\mathcal U$, and put
\[
 R_\kappa=\mathbb Q^\kappa/J_{<\kappa},\qquad
 J_{<\kappa}=\{r:|\operatorname{supp}(r)|<\kappa\}.
                                                               \tag{H5.9.4}
\]
The denominator is an ideal. Under the no-measurable hypothesis,
$R_\kappa$ has no nonzero finite-dimensional unital quotient.
For otherwise a maximal ideal in that quotient gives a unital
map $h:R_\kappa\to K$, where $K/\mathbb Q$ is a finite field
extension. In fact $K=\mathbb Q$. To see this, represent any
$r\in R_\kappa$ by rational coordinates $r_i$. If the minimal
polynomial of $h(r)$ were irreducible of degree greater than one,
it would have no rational zero. Its value on $(r_i)$ would then
be a unit in $\mathbb Q^\kappa$, with its pointwise reciprocal
as inverse, whereas its image in $K$ would be zero. This is
impossible. Thus every $h(r)$ is rational; the map is onto its
residue field, which proves $K=\mathbb Q$.

For $A\subseteq\kappa$ write $e_A$ for its characteristic
idempotent. The composite $\mathbb Q^\kappa\to R_\kappa
\xrightarrow{h}\mathbb Q$ defines an ultrafilter
\[
                  \mathcal D=\{A:h(e_A)=1\}.
\]
More explicitly, Boolean identities for the $e_A$ prove the
ultrafilter axioms, and $h$ kills all $e_A$ with $|A|<\kappa$;
thus $\mathcal D$ is uniform and nonprincipal. For any rational
function $r=(r_i)$, put $a=h(r)$. The set
$\{i:r_i=a\}$ belongs to $\mathcal D$. If it did not, its
complement would belong to $\mathcal D$; defining the reciprocal
of $r-a$ on that complement and zero elsewhere would give
$(r-a)s=e_{\{r_i\ne a\}}$, with images $0=1$.
Apply this to a function taking distinct integer values on the
parts of any countable partition of $\kappa$. Some part must
belong to $\mathcal D$. This is equivalent to countable
completeness of $\mathcal D$.

The general ultrafilter implication is also stated in
\citet[Theorem 1.1]{EH5BSZ2026}. We include its short proof
to ensure that the resulting measurable cardinal lies inside
$\mathcal U$. Choose the least cardinal $\mu\leq\kappa$ for
which $\kappa$ has a partition into $\mu$ parts none in
$\mathcal D$; the singleton partition proves existence.
Countable completeness makes $\mu$ uncountable. Minimality
makes $\mathcal D$ closed under intersections of fewer than
$\mu$ members: a contrary family, partitioned according to
the first failed member and its total intersection, would
give a smaller partition with no part in $\mathcal D$.
The induced ultrafilter on the set of $\mu$ partition classes
is a nonprincipal $\mu$-complete ultrafilter. Hence $\mu$
is measurable, a contradiction. We have proved the assertion
about finite-dimensional quotients of $R_\kappa$ without
any compactness assumption.

The detection lemma now implies that the quotient
\[
 q_\kappa:\mathcal H(R_\kappa)\longrightarrow
                \mathcal H(R_\kappa)/C(R_\kappa)=R_\kappa^2
                                                               \tag{H5.9.5}
\]
belongs to the group equivalence class defining
$W_{\mathrm{unip}}$: every allowed representation kills its
kernel, and the factorization through a quotient is unique.
The source group, however, is not local for this class.
Its identity map cannot factor through $q_\kappa$, because
$[(1,0,0),(0,1,0)]=(0,0,1)\ne1$.

Suppose now that a set $S\in\mathcal U$ of functors generated
$W_{\mathrm{unip}}$. Choose an uncountable regular
$\kappa\in\mathcal U$ larger than the numbers of simplices
in the nerves of all their sources and targets.
Products of copies of $B\mathcal H(\mathbb Q)$ are $S$-local,
since this allowed test group is $W_{\mathrm{unip}}$-local.
There is a canonical group isomorphism
\[
 \mathcal H(R_\kappa)\cong
  \colim_{A\in[\kappa]^{<\kappa}}
              \mathcal H(\mathbb Q)^{\kappa\setminus A},
                                                               \tag{H5.9.6}
\]
where the transition maps forget coordinates. Its indexing
poset is $\kappa$-filtered by regularity. Indeed both sides
identify coordinate families which agree outside a set of
cardinality less than $\kappa$, and the multiplication
(H5.9.1) is coordinatewise. Nerves of groups commute with
filtered colimits, and those nerves are Kan complexes.
Mapping out of any simplicial set with fewer than $\kappa$
simplices commutes with $\kappa$-filtered colimits: in each
simplicial degree a map and all its required equalities
involve fewer than $\kappa$ simplices. Filtered colimits
preserve weak equivalences of Kan complexes, as seen by
representing a sphere, disk, or homotopy in one stage.
It follows that $B\mathcal H(R_\kappa)$ is $S$-local.
If $S$ generated the stated localizer, it would be
$W_{\mathrm{unip}}$-local, contradicting (H5.9.5).
This proves nonaccessibility.

Under Vop\v enka's principle the positive assertion is
\citet[Scholium 4.2.17]{EH5Cisinski2006}, whose exact
conclusion is that every fundamental localizer is accessible.
Its application is legitimate because H5.1 and LM8 B.6
establish the relative Theorem A, not merely absolute
asphericity closure. Given a small generating set, bound
the sizes of its source and target fundamental groups by
a regular $\kappa$. Since $W_1$ is already generated by
$S^2\to *$, replacing each generator by its component
fundamental-group homomorphisms shows that the set described
in the statement contains a generating set and is contained
in $W_{\mathrm{unip}}$. It is therefore itself a generating set.
The two implications distinguish the fixed finite-dimensional
test definition from the different localization at all uniquely
divisible nilpotent groups; those targets cannot be substituted.
\end{proof}
